\documentclass[11pt]{article}
\usepackage{amsmath,amssymb,amsthm}
\newtheorem{theorem}{Theorem}
\title{Positive Pair Compatibility Fails on the Hierarchical Pole Grid}
\author{}\date{}
\begin{document}\maketitle
\paragraph{Status.}
This is a complete finite-certificate obstruction, verified by exact
rational arithmetic and guarded formal power series. It is not Lean
checked. It refutes a positive-residue induction invariant, not the
reverse compiler or either requested global asymptotic statement.
Signed residues, target-dependent poles, and other constructions remain
outside its scope. The original model and its definitions are unchanged.

Write $s=s_0+2s_1+4s_2$ and put
\[
 p_s=s_0\epsilon+s_1\epsilon^2+s_2\epsilon^4,
 \qquad \delta=\epsilon^8,\qquad t_s=1-p_s.
\]
The sixteen poles are $p_s,p_s+\delta$, for $0\leq s<8$.
For $0<\epsilon\leq1/4$ they are all below $1/2$, so the
eight queries $t_s$ are strictly exterior. Given a word
$w\in\{-1,1\}^8$ and a shared residue vector $c\in\mathbb R^{16}$,
define
\[
 F_w(t)=\sum_{s=0}^7w_s\left(
 \frac{c_s}{t-p_s}-\frac{c_{s+8}}{t-p_s-\delta}\right).
\]
\begin{theorem}
Let
\[
 \sigma=(1,1,1,1,-1,-1,-1,-1),\qquad
 \tau=(1,-1,1,1,-1,-1,1,1).
\]
For every sufficiently small $\epsilon>0$, there are no
nonnegative shared residues $c_s\geq0$ and separate biases
$\beta_\sigma,\beta_\tau$ satisfying
\[
 w_s\bigl(F_w(t_s)+\beta_w\bigr)>0
 \quad\text{for all }s\text{ and }w\in\{\sigma,\tau\}.
\]
The same nonexistence holds at $\epsilon=1/100$.
\end{theorem}
\begin{proof}
For each word, elimination of its bias gives the necessary inequalities
$F_w(t_a)>F_w(t_b)$ for every $w_a=1,w_b=-1$. Dividing by
$|p_a-p_b|>0$ gives a rational row $R_{w,a,b}$ with
\[
 (R_{w,a,b})_s=
 \frac{\operatorname{sgn}(p_a-p_b)w_s}
 {(1-p_a-p_s)(1-p_b-p_s)},\qquad
 (R_{w,a,b})_{s+8}=
 \frac{-\operatorname{sgn}(p_a-p_b)w_s}
 {(1-p_a-p_s-\delta)(1-p_b-p_s-\delta)}.
\]
For $0<\epsilon\leq1/4$, the sign in this formula is determined
by the least binary position where $a$ and $b$ differ; the later
powers have smaller total magnitude. Thus the rows are fixed rational
functions of $\epsilon$ throughout that interval.

Take the following eleven comparison rows, in the displayed order:
\[
 \begin{split}
 &(\sigma,0,5),(\sigma,0,7),(\sigma,1,4),(\sigma,1,6),
   (\sigma,2,5),(\sigma,2,6),(\sigma,3,4),\\
 &(\tau,0,1),(\tau,0,4),(\tau,2,4),(\tau,3,4).
 \end{split}
\]
Append the six coordinate rows
$e_0^\top,e_2^\top,e_4^\top,e_{13}^\top,e_{14}^\top,e_{15}^\top$.
Call this $17$ by $16$ matrix $A(\epsilon)$. Define
\[
 M(\epsilon)=\begin{pmatrix} A(\epsilon)^\top\\
                      1&\cdots&1\end{pmatrix},\qquad
 \lambda(\epsilon)=M(\epsilon)^{-1}e_{17}.
\]
Here $e_{17}$ denotes the last coordinate vector in $\mathbb R^{17}$.
Exact elimination verifies that $M$ is invertible over
$\mathbb Q(\epsilon)$ and that the seventeen leading terms of
$\lambda$, in this row order, are
\[
 \begin{split}
 (&3\epsilon/2,1/4,3\epsilon/2,1/4,1/4,3\epsilon^4,1/4,\\
 &\epsilon^{11}/8,\epsilon^{11}/8,
   \epsilon^{11}/8,\epsilon^{11}/8,\\
 &\epsilon^{28}/4,\epsilon^{26}/4,\epsilon^{28}/4,
   \epsilon^{26}/4,\epsilon^{25}/4,\epsilon^{26}/4).
 \end{split}
\]
Each leading coefficient is strictly positive. Consequently all
entries of $\lambda(\epsilon)$ are positive for every sufficiently
small positive $\epsilon$. Its defining equations give
$\lambda^\top A=0$ and $\sum_i\lambda_i=1$ exactly.

For reproducibility, \texttt{positive-paired-root-h16-check.py}
regenerates all rows and verifies these leading terms using exact
rational power series at precision $512$. It always pivots at the
smallest detected valuation, guards every truncated zero before
division, and propagates every multiplier, including truncated zeros.
The pivot valuations are
\[
 (0,0,0,0,0,0,0,0,0,2,2,6,6,8,15,17,23).
\]
It scales the right side by $\epsilon^{170}$ during the calculation
so that intermediate series have nonnegative valuation. Every
displayed leading term is determined below the tracked precision.
These guarded operations agree with elimination over the rational
function field, so the resulting leading signs are exact, rather
than numerical approximations. The saved rational vector at
$\epsilon=1/100$ is checked separately without using a solver.

Under the proposed strict root inequalities, the first eleven
entries of $Ac$ would be positive, while its final six entries would
be nonnegative. Thus $\lambda^\top Ac>0$, contradicting
$\lambda^\top A=0$. This proves both assertions.
\end{proof}

\paragraph{Consequence for the proposed induction.}
These are the root profiles of two sixteen-wire butterfly targets.
Lower-layer routing preserves the highest input bit, so the root
profile after routing is precisely $(w_s,-w_s)$, independently of
the lower switches. In particular no two full targets with the
displayed root words can share a nonnegative residue vector at
sufficiently small scale, even when every block is allowed its own
target-dependent bias. This refutes a universal positive child-cone
pairing invariant at this fixed size.

The dual argument permits residues of arbitrary size and arbitrary
dependence on $\epsilon$. Diverging positive residue germs therefore
do not evade it. The coordinate rows are essential to its scope:
the result makes no assertion for shared signed residues, and makes
no assertion against separately chosen residues for the two targets.
\end{document}
